% ID: SEL-D0015-BD-E08
% ESTADO: cotejada. Segunda pasada visual completa y borrador compilado revisados.
% Fuente: Selectividad/matematicas_aplicadas/matematicas_aplicadas_P_C.pdf
% SHA-256: 463bf63d4697a18a0b9efc9287e7b0a593c803bab98a3cfab0e260067a0dd02c
% Página: 2; bloque D; ejercicio original 8.

La renta anual de los hogares andaluces, en miles de euros, se distribuye
según una ley Normal con desviación típica 5 y media desconocida $\mu$.
\begin{enumerate}[label=\alph*)]
  \item \textbf{(1 punto)} Si se desea que en el 99\% de las posibles
  muestras del mismo tamaño, elegidas de entre los hogares andaluces, la
  media muestral no difiera de la renta media anual poblacional de dichos
  hogares en más de una unidad, \textquestiondown{}cuál debe ser el tamaño
  mínimo de las muestras?
  \item \textbf{(0.5 puntos)} Si se consideran muestras de hogares andaluces
  de tamaño 100, \textquestiondown{}qué distribución de probabilidad sigue
  la variable aleatoria ``Renta media anual muestral''?
  \item \textbf{(1 punto)} Suponiendo que la renta media anual poblacional
  de los hogares andaluces es $\mu=24$, \textquestiondown{}cuál es la
  probabilidad de que en una muestra de tamaño 100 la renta media anual
  muestral sea superior a 25?
\end{enumerate}
